\section{The ordered microscopic tail and the original pointwise law}
\label{sec:v83-pointwise}
The preceding estimates give complete source and surface controls in
specified local norms. The original pointwise target is stronger.
We now state the remaining implication on that target, and separate
the microscopic tail before any interchange of limits could occur.

For each exact label retain the actual source $v^\eta$. For $r>0$,
put $\ell_j=2^{-j}r$ and define
\begin{equation}\label{eq:v83-dyadic-tail}
 T_{\eta,r,N}(t)=\sum_{j=N}^{\infty}\ell_j^{-\kappa_\eta}
                  \int_{\ell_j/2}^{\ell_j}v^\eta(t+s)\,ds.
\end{equation}
The integration variable $s$ is an offset from the starting roof.
Equivalently one integrates $v^\eta(u)\,du$ over
$(t+\ell_j/2,t+\ell_j]$; there is no second translation of the roof.
All terms and the entire tail are formed at finite collision count.

\begin{lemma}[An ordered decomposition, including the full tail]
\label{lem:v83-ordered-tail}
Pointwise as extended nonnegative quantities,
\begin{equation}\label{eq:v83-potential-dyadic}
 T_{\eta,r,0}(t)\le\mathsf P_{\eta,r}v^\eta(t)
                  \le2^{\kappa_\eta}T_{\eta,r,0}(t).
\end{equation}
Define
\begin{align}
 \mathfrak T_{\eta,M}(r)
  &=\lim_{N\to\infty}\limsup_m\sup_{R,n,k}m^2
       \esssup_{t\in\mathcal C_{m,R,n,k}(M)}T_{\eta,r,N}(t),
                                                 \label{eq:v83-tail-height}\\
 \mathfrak P_{\eta,M}(r)
  &=\limsup_m\sup_{R,n,k}m^2
       \esssup_{t\in\mathcal C_{m,R,n,k}(M)}
                                     \mathsf P_{\eta,r}v^\eta(t).
                                                 \label{eq:v83-potential-height}
\end{align}
The first limit exists monotonically in $[0,\infty]$, and
\begin{equation}\label{eq:v83-tail-bound}
 \mathfrak T_{\eta,M}(r)\le\mathfrak P_{\eta,M}(r)
       \le C_\eta r^{(1+\eta)/2}
                            +2^{\kappa_\eta}\mathfrak T_{\eta,M}(r).
\end{equation}
\end{lemma}
\begin{proof}
On $(\ell_j/2,\ell_j]$, the kernel lies between
$\ell_j^{-\kappa_\eta}$ and $2^{\kappa_\eta}\ell_j^{-\kappa_\eta}$.
Tonelli proves \eqref{eq:v83-potential-dyadic}. For any fixed $N$,
the first $N$ interval lengths are fixed before the collision limit.
Input~\ref{input:v83-local} bounds the normalized height of their
integrals by $C_\eta\ell_j$. Hence their complete finite sum is at
most $C_\eta\sum_{j<N}\ell_j^{1-\kappa_\eta}\le
C_\eta r^{1-\kappa_\eta}$. Keep the whole remaining nonnegative
tail inside the supremum and collision limsup. Only then let
$N\to\infty$. Monotonicity of these tails proves existence and
\eqref{eq:v83-tail-bound}. This argument does not use the false
reverse of Fatou's inequality for an infinite sum.
\end{proof}

Let the full, untruncated surface height be
\begin{equation}\label{eq:v83-surface-height}
 \mathfrak F_M^{\rm all}(r)=\limsup_m\sup_{R,n,k}m^2
       \esssup_{t\in\mathcal C_{m,R,n,k}(M)}
                                    \overline{\mathcal F}_r(t).
\end{equation}
Its source is precisely \eqref{eq:v83-untruncated-flux}.

\begin{theorem}[Pointwise error with no initial-cutoff interchange]
\label{thm:v83-endpoint}
For each fixed $0<\eta<1/2$, all sufficiently large fixed $B$ and
all previously fixed $0<r<r_0$, the original pointwise error satisfies
\begin{align}
 \mathcal E_M\le{}&C_MB^{-1/192}+C\mathfrak F_M^{\rm all}(r)
                         +C_\eta\mathfrak T_{\eta,M}(r)\notag\\
 &+C_\eta r^{(1+\eta)/2}
             +C_\Theta(1+\varepsilon(B)^{-1})\sqrt r.
                                                 \label{eq:v83-endpoint}
\end{align}
Thus the original assertion \eqref{eq:v76-target} follows if the
full surface height and the microscopic tail in this display tend
to zero as $r\downarrow0$. This sufficient condition is not asserted
as proved by the local-norm estimates above.
\end{theorem}
\begin{proof}
Use the complete-source shift inequality of
Input~\ref{input:v83-geometric} and replace its sum of true-exit
and optical losses, at the source level, by the average of
\eqref{eq:v83-source-upper}. The positive scalar remainder and the
fixed-width bad-source bound then give
\[
 \mathcal E_M\le C_MB^{-1/192}+C\mathfrak F_M^{\rm all}(r)
   +C_\eta\mathfrak P_{\eta,M}(r)
   +C_\Theta(1+\varepsilon(B)^{-1})\sqrt r.
\]
Apply \eqref{eq:v83-tail-bound}. Every source restriction has already
been removed at fixed count, using actual fractional integrability;
no first inverse-gradient cutoff is passed through a collision limit.
To obtain the conditional conclusion, choose
$r(B)=c\varepsilon(B)^\gamma$ with fixed $c>0$ and $\gamma>2$.
For each $B$ this is a fixed positive width. First take the collision
limit, next enlarge auxiliary spectral bands, and take $B\to\infty$
last. Every displayed elementary error then vanishes, and the two
stated conditions pay the remaining terms. The arithmetic reference
and all zero classes are unchanged.
\end{proof}

\begin{proposition}[A quantitative microscopic envelope]
\label{prop:v83-envelope}
Suppose, as an additional estimate on the actual source, that for
all sufficiently large $m$, all parameters and almost every central
$t$,
\begin{equation}\label{eq:v83-envelope}
 m^2\int_0^s v^\eta(t+u)\,du\le H(s)\quad(0<s<r_0),\qquad
 \int_0^{r_0}H(s)s^{-1-\kappa_\eta}\,ds<\infty,
\end{equation}
with $H$ nonnegative and nondecreasing. The exceptional set may be
chosen independently of $s$ by first using rational $s$ and monotone
limits. Then $\mathfrak T_{\eta,M}(r)=0$ for every $r<r_0/2$.
In particular a uniform bound $H(s)\le A s^{\kappa_\eta+\delta}$,
$\delta>0$, gives the finite-count tail bound
\begin{equation}\label{eq:v83-envelope-tail}
 m^2 T_{\eta,r,N}(t)
        \le \frac{A r^\delta2^{-N\delta}}{1-2^{-\delta}}.
\end{equation}
The envelope is a uniform finite-count hypothesis, not a list of
separate fixed-scale collision limsups.
\end{proposition}
\begin{proof}
Each normalized interval integral is at most $H(\ell_j)$.
Moreover, monotonicity gives
\[
 H(\ell_j)\ell_j^{-\kappa_\eta}
 \le C_\eta\int_{\ell_j}^{2\ell_j}H(s)s^{-1-\kappa_\eta}\,ds.
\]
The right intervals have disjoint interiors. Their tails tend to
zero by the assumed integrability, uniformly before the collision
limit. For the power envelope sum the geometric series directly.
\end{proof}

The fractional source lowers the sufficient envelope exponent from
$1/2$ to $\kappa_\eta=(1-\eta)/2$, which can be any fixed number
strictly between $1/4$ and $1/2$. The endpoint $1/4$ is not included.
The complete surface mass bound \eqref{eq:v83-surface-local-mass}
and the weak-$L^{Q_\eta}$ optical bound do not establish
\eqref{eq:v83-envelope} or \eqref{eq:v83-surface-height}. A uniformly
small integral, even with these additional local controls, must not
be substituted for an essential roof-height estimate. The original
pointwise law, and scalar-first normalization of its path and
likelihood consequences, remain exactly the organizing objective.
